\chapter{编辑器与终端工作流}
\label{ch:editor-workflow}
\chapterlead{日常效率取决于三件事：编辑器的操作是否成了肌肉记忆，终端能否同时管理多个任务，shell 的补全与历史是否好用。本章以 Emacs 为主，在终端里使用；vim 作为服务器上的兜底。配置一律交给 Nix 与 Home Manager 生成。}

\section{起步：在终端里启动 Emacs}

\begin{shellsession}
$ emacs -nw
$ emacs -nw notes.txt
\end{shellsession}

\texttt{-nw} 是 \texttt{-{}-no-window-system} 的缩写，要求 Emacs 不开图形窗口，直接在终端里绘制界面。加上它，SSH 登录的服务器与本机的使用方式完全一致；不加时 Emacs 会先尝试连接显示服务，失败才回退到终端。

启动后屏幕下方是回显区，显示按键与提示。光标所在位置即编辑位置，直接输入文字。

十个动作覆盖日常编辑。

\begin{booktable}{\textwidth}{L{0.92} L{1.02} J{1.06}}
\toprule
按键 & 作用 & 说明 \\
\midrule
\endhead
\texttt{Ctrl+x} \texttt{Ctrl+f} & 打开文件 & 输入路径，不存在则新建 \\
\texttt{Ctrl+x} \texttt{Ctrl+s} & 保存 & 不退出，可继续编辑 \\
\texttt{Ctrl+x} \texttt{Ctrl+c} & 退出 & 有未保存文件时会逐个询问 \\
\texttt{Ctrl+s} & 向前搜索 & 边输入边跳转，再按继续找下一个 \\
\texttt{Ctrl+r} & 反向搜索 & 与 \texttt{Ctrl+s} 方向相反 \\
\texttt{Ctrl+g} & 取消 & 放弃当前命令、退出回显区 \\
\texttt{Alt+x} & 按名执行命令 & 输入命令名，Tab 可补全 \\
\texttt{Ctrl+space} & 开始选区 & 移动光标即选中；被输入法占用时改用 \texttt{Ctrl+@} \\
\texttt{Ctrl+w}、\texttt{Alt+w} & 剪切、复制 & 对选区生效 \\
\texttt{Ctrl+y} & 粘贴 & 贴回最近剪切或复制的内容 \\
\bottomrule
\end{booktable}

\begin{notebox}[出问题先按 Ctrl+g]
Emacs 的按键是多层组合，误触后容易停在半截状态：回显区提示 \texttt{C-x}，或在询问一个没读懂的选项。此时按 \texttt{Ctrl+g} 回到正常状态，不会丢失文件内容。终端里 \texttt{Ctrl+s} 若被流控截住，按 \texttt{Ctrl+q} 恢复，原因见第\ref{ch:first-terminal}章。
\end{notebox}

\subsection{以 Emacs 为主的理由}

第一，键位与命令行同源。\texttt{Ctrl+a}、\texttt{Ctrl+e}、\texttt{Ctrl+k}、\texttt{Ctrl+y} 在提示符下与 Emacs 里含义相同，一套肌肉记忆覆盖两处。readline 的键位就是照 Emacs 定的，第\ref{ch:first-terminal}章已经用过其中一部分。

第二，终端里运行，服务器与本机一致。\texttt{emacs -nw} 不依赖图形会话，通过 SSH 使用同一套操作与配置。

第三，配置是代码。Emacs 的配置与扩展用 Emacs Lisp 写成，可以进版本控制、可以在 Nix 表达式里生成，与本书其余部分的做法一致。

Emacs 的代价同样存在：默认启动比 vim 慢，界面朴素，扩展生态的学习曲线不低。终端版（\texttt{emacs-nox}）已经去掉了图形依赖，剩下的启动时间与配置成本无法完全消除。

vim 仍然是兜底。任何一台 Linux 上都可能只有 \texttt{vi}，包括救援模式与最小化安装的容器；服务器上改一个配置文件，用 vim 的六个动作最快。两条路线不冲突：日常编辑用 Emacs，改动单个配置文件用 vim。

\section{用 Nix 安装 Emacs}

系统级安装只需一行。

\begin{nixcode}
environment.systemPackages = with pkgs; [ emacs-nox ];
\end{nixcode}

\texttt{emacs-nox} 是不带图形界面支持的构建版本：没有 GTK 与 X11 依赖，安装体积小，在任何终端里都能启动。它与 \texttt{emacs} 的版本相同，功能差别仅在图形部分。日常在终端里使用，装这个即可。

\begin{shellsession}
$ emacs --version | head -1
GNU Emacs 31.1
\end{shellsession}

版本号随 nixpkgs 走，写在配置里的系统代际连同 Emacs 一起回滚。

\section{用 Home Manager 配置 Emacs}

编辑器的配置属于用户，放在 Home Manager 里。

\begin{nixcode}
programs.emacs = {
  enable = true;
  package = pkgs.emacs-nox;
  extraPackages = epkgs: [
    epkgs.magit
    epkgs.markdown-mode
  ];
  extraConfig = ''
    (setq make-backup-files nil)
    (setq-default indent-tabs-mode nil)
    (global-display-line-numbers-mode 1)
    (setq inhibit-compacting-font-caches t)
  '';
};
\end{nixcode}

四项设置分别对应：不生成 \texttt{foo\textasciitilde} 备份文件、用空格代替制表符、显示行号、避免中文字体缓存反复重建。\texttt{extraPackages} 选择要预装的扩展；上例是 Git 前端与 Markdown 模式。

\texttt{extraConfig} 里的内容写进 site 级的 \texttt{default.el}，不是 \texttt{\textasciitilde/.emacs.d/init.el}。这样设计是为了不与个人的 Emacs 配置框架冲突：若以后引入 Doom 或 Purcell 的配置集，\texttt{default.el} 仍然生效，Nix 算出的路径也能被它们使用。

\begin{shellsession}
$ emacs --batch --eval '(princ (locate-library "default"))'
/nix/store/...-default-0.1.0/share/emacs/site-lisp/elpa/default-0.1.0/default.el
\end{shellsession}

输出落在 \texttt{/nix/store} 下，说明配置由 Nix 生成。修改 \texttt{extraConfig} 后执行一次 \texttt{home-manager switch}，重新打开 Emacs 即生效。

\begin{notebox}[图形界面与终端共用一份配置]
需要图形界面时，把 \texttt{package} 换成 \texttt{pkgs.emacs}（含 GTK）。命令行加 \texttt{-nw} 仍进终端模式，配置与扩展是同一套，无需分别维护。
\end{notebox}

\section{vim 的六个动作}

服务器上没有 Emacs、或只需要改一个文件时，用 vim。

\begin{shellsession}
$ vim --version | head -2
$ vim notes.txt
\end{shellsession}

vim 分模式：普通模式下按键是命令，插入模式下按键是文字。六个动作覆盖大部分改配置的场景。

\begin{booktable}{\textwidth}{J{0.720} J{1.280}}
\toprule
按键 & 作用 \\
\midrule
\endhead
\texttt{i} & 进入插入模式，光标前插入 \\
\texttt{Esc} & 回到普通模式 \\
\texttt{:w} & 保存 \\
\texttt{:q} & 退出；有未保存改动时用 \texttt{:q!} 放弃 \\
\texttt{dd} & 删除当前行 \\
\texttt{yy} 与 \texttt{p} & 复制当前行，粘贴到下一行 \\
\bottomrule
\end{booktable}

再加两个动作：\texttt{/关键词} 搜索，\texttt{n} 跳到下一个匹配；\texttt{u} 撤销。\texttt{:\%s/old/new/g} 全文替换，加 \texttt{c} 结尾则每次替换前确认。

\section{Neovim 与最小配置}

Neovim 是 vim 的重写版本，配置语言是 Lua，生态更活跃。若选 vim 路线，用 Home Manager 生成最小配置。

\begin{nixcode}
programs.neovim = {
  enable = true;
  defaultEditor = true;
  viAlias = true;

  extraLuaConfig = ''
    vim.opt.number = true
    vim.opt.relativenumber = true
    vim.opt.expandtab = true
    vim.opt.shiftwidth = 2
    vim.opt.ignorecase = true
    vim.opt.smartcase = true
  '';
};
\end{nixcode}

六行设置分别对应：显示行号、相对行号（便于按行数跳转）、用空格代替制表符、缩进宽度、搜索忽略大小写、但含大写时区分。

需要语法高亮与语言服务时，可以引入现成的发行版，例如 LazyVim，或用 \texttt{nixvim} 把配置写成 Nix 表达式。引入前先确认需要哪些语言，避免一次装上几十个用不到的插件。

\begin{notebox}[先学操作，再学配置]
编辑器的配置是无底洞。前两周只用默认配置，把上面的动作练到不用想，再考虑扩展。配置的价值建立在操作熟练之上，顺序颠倒会花掉大量时间在调整键位上。
\end{notebox}

\section{终端复用：tmux 与 zellij}

终端复用器让一个窗口里管理多个会话，断开 SSH 后会话仍在。

\begin{shellsession}
$ tmux new -s work      # 新建会话
$ tmux ls               # 列出会话
$ tmux attach -t work   # 重新接入
\end{shellsession}

tmux 的默认前缀键是 \texttt{Ctrl+b}，按下后接命令：\texttt{c} 新建窗口，\texttt{n} 下一个窗口，\texttt{"} 水平分屏，\texttt{\%} 垂直分屏，\texttt{d} 断开。zellij 提供类似的界面，底部常驻键位提示，不必记前缀键，代价是占用一行屏幕。

在服务器上，复用器的价值最大：网络中断、笔记本合盖、SSH 客户端崩溃，都不会终止会话中正在运行的编译或部署。在 tmux 里运行 \texttt{emacs -nw}，同样不受断线影响。

\section{shell 的选择与补全}

\begin{nixcode}
programs.zsh = {
  enable = true;
  enableCompletion = true;
  autosuggestion.enable = true;
  syntaxHighlighting.enable = true;
  initExtra = ''
    bindkey -e
    export EDITOR=emacs
  '';
};

programs.starship.enable = true;
\end{nixcode}

三项能力决定交互体验：补全（按 Tab 补齐参数与路径）、历史搜索（\texttt{Ctrl+r}）、语法高亮（命令存在与否一眼可见）。\texttt{bindkey -e} 让 zsh 使用 Emacs 键位，与 readline 和 Emacs 保持一致。\texttt{EDITOR=emacs} 之后，第\ref{ch:first-terminal}章讲的 \texttt{Ctrl+x} \texttt{Ctrl+e} 会用 Emacs 编辑当前命令行。

脚本仍用 bash 编写。原因在第\ref{ch:first-terminal}章说过：服务器上一定有 bash，脚本要选最通用的实现。

\section{把配置纳入版本控制}

编辑器、终端、shell 与行编辑的配置都应当由 Home Manager 生成，进入同一个仓库。

\begin{nixcode}
# home.nix 片段
home.file.".config/kitty/kitty.conf".text = ''
  font_family      Recursive Mono Casual
  font_size        12
  background_opacity 0.95
  map ctrl+shift+c copy_to_clipboard
'';

home.file.".inputrc".text = ''
  set editing-mode emacs
  set completion-ignore-case on
  set bell-style none
'';
\end{nixcode}

\texttt{inputrc} 只对使用 readline 的程序生效，例如 bash 与 \texttt{psql}；zsh 的键位由前面配置的 \texttt{bindkey -e} 决定，不读这份文件。

\texttt{home.file} 适合写小段静态配置，\texttt{programs.*} 适合程序有专用模块的情况。Emacs 用 \texttt{programs.emacs}，不要用 \texttt{home.file} 另写一份 \texttt{init.el}，否则与 \texttt{extraConfig} 各管一半，排查困难。两者生成的路径都在 store 里，\texttt{\$HOME} 下是指向它的链接，源文件随仓库走。

\begin{shellsession}
$ ls -l ~/.config/kitty/kitty.conf ~/.inputrc
$ git -C ~/nixos-config log --oneline -3
\end{shellsession}

第一条确认文件来源，第二条确认改动有历史。这两条命令是“配置可追溯”的最小验证。

\section{练习}

\begin{enumerate}
  \item 在终端里启动 \texttt{emacs -nw}，用十个动作新建一个文件、写入三行、保存退出，再用 \texttt{Ctrl+x} \texttt{Ctrl+f} 打开它。
  \item 在 Emacs 里用 \texttt{Ctrl+space} 选中一段文字（被输入法占用时用 \texttt{Ctrl+@}），用 \texttt{Ctrl+w} 剪切，移到文末用 \texttt{Ctrl+y} 贴回。
  \item 用 \texttt{emacs --batch --eval} 执行 \texttt{(princ (locate-library "default"))}，找出配置文件的路径，说明它为什么在 store 里。
  \item 用 vim 完成同样的事：打开文件、插入三行、保存退出；再用 \texttt{u} 撤销一次改动。
  \item 设 \texttt{EDITOR=emacs}，在提示符下按 \texttt{Ctrl+x} \texttt{Ctrl+e} 写一条三行的循环命令，保存退出后确认它被执行。
  \item 用 tmux 建立两个窗口，一个跑 \texttt{emacs -nw}，断开后重新接回，确认编辑内容仍在。
  \item 用 \texttt{home.file} 写一份 \texttt{\textasciitilde/.inputrc}，切换后确认 \texttt{\$HOME} 下是符号链接，并在 bash 里用 \texttt{bind -f} 载入它。
\end{enumerate}

\begin{checkpoint}
本章结束时，应能：
\begin{itemize}
  \item 在终端里启动并退出 Emacs，用十个动作完成打开、编辑、保存与搜索；
  \item 说明终端路线（\texttt{emacs-nox} 加 \texttt{-nw}）在服务器上的好处；
  \item 用 \texttt{environment.\allowbreak{}systemPackages} 与 \texttt{programs.emacs} 分别完成安装与配置；
  \item 说明 \texttt{extraConfig} 写入 \texttt{default.el} 而非 \texttt{init.el} 的原因；
  \item 用 vim 的六个动作在无 Emacs 的机器上改配置文件；
  \item 用 tmux 或 zellij 管理会话，说明它解决了什么问题；
  \item 配置 shell 的补全、历史搜索与语法高亮，并让 zsh 使用 Emacs 键位；
  \item 把编辑器、终端与行编辑的配置纳入 Home Manager 与版本控制。
\end{itemize}
\end{checkpoint}
